\chapter{曲风与配器：八种真实写法}
\label{ch:genres}

\classquote{易穷则变，变则通，通则久。}{《周易·系辞下》}{走到尽头就要变，变了才走得通，通了才能长久。曲风从来不是死标签，它一直在借鉴与变形。曲风不是贴在声音上的标签，而是一套“用什么乐器、怎样分行、把重音放在哪里”的选择。学会一套选择，比背十个 Preset 更有用。}

\section{为什么把曲风独立成章}
\label{sec:genres-why}

前九章按“功能层”和“乐器族”拆开讲：节奏、和声、旋律、弦乐、木管铜管、电子、人声。这样适合建立概念，但真正写歌时，这些选择是同时发生的：一首摇滚歌的\term{节奏组}{rhythm-section}会决定吉他怎么写，一首电子舞曲的\term{低频层}{sub-layer}会决定贝斯和底鼓怎样让位。

本册把曲风内容独立成第 11 章，理由有三条：

\begin{enumerate}
\item \textbf{不重复}：如果把八种曲风拆进前十章，同一套“编制—律动—和声—结构”会在不同章节反复出现，读者很难横向比较。
\item \textbf{可查阅}：写摇滚时不需要翻完爵士和电子两节；本章按风格分节，每节都能独立照着做。
\item \textbf{不替代基础}：曲风是选择，不是规则。先读第 1—5 章的四层角色、段落、律动、和声与旋律；本章把这些选择装进具体的风格骨架。
\end{enumerate}

\section{配器四把尺}
\label{sec:orchestration-rulers}

在进入具体曲风之前，先把四个判断工具磨快。它们适用于所有风格，也是“为什么这样配”的共同语言。

\subsection{音域与移调}
\label{sec:range-transposition}

\term{音域}{range}是一件乐器或一个人声能发出的最低到最高音；\term{适奏音区}{tessitura}是它最舒服、最常待的那一段。写配器时，先问“这个角色应该待在哪一段”，再问“这件乐器能不能到那里”。

\begin{longtable}{p{0.22\textwidth}p{0.30\textwidth}p{0.38\textwidth}}
\toprule
乐器/声部 & 常用音域（记谱） & 移调与备注 \\
\midrule
\endfirsthead
\toprule
乐器/声部 & 常用音域（记谱） & 移调与备注 \\
\midrule
\endhead
小提琴 & G3–E7 & 不移调；高音区明亮，低音区不适合长期铺底 \\
中提琴 & C3–A6 & 不移调；中音厚，适合内声部 \\
大提琴 & C2–C6 & 不移调；低音锚与歌唱性旋律都能做 \\
低音提琴 & E1–G4 & 实际发音低八度；记谱常按实际音高写 \\
长笛 & C4–C7 & 不移调；高音穿透，低音区气息声明显 \\
双簧管 & B$\flat$3–F6 & 不移调；音色集中，适合独奏句 \\
单簧管（B$\flat$） & 记谱 E3–C7 & 实际发音低大二度；写谱前先确认音色库按记谱还是实际音高 \\
巴松 & B$\flat$1–F5 & 不移调；低音木管，不适合频繁跳高 \\
小号（B$\flat$） & 记谱 F$\sharp$3–C6 & 实际发音低大二度；高音区要留体力 \\
圆号（F） & 记谱 C3–C6 & 实际发音低纯四度；常写中低音区 \\
长号 & E2–B$\flat$4 & 不移调；低音铜管的骨干 \\
人声 & 见下 & 按角色选，不按极限选 \\
\bottomrule
\end{longtable}

人声的经验起点：女高音 C4–A5、女中音 G3–E5、男高音 C3–A4、男低音 E2–E4。真正写歌时，把主唱最舒服的两三个音当“家”，副歌再向上借三到五度；不要为了高音把整首歌都推到极限。

\term{移调}{transposition}有两层：一是记谱移调（B$\flat$ 单簧管、B$\flat$ 小号、F 圆号等），二是为了音区或音色把整段旋律挪动。夜半的 MIDI 默认按实际音高（concert pitch）发声；如果你加载的是按记谱音高制作的音色，必须按音色说明换算，不能把乐谱上的 written pitch 直接输入。最稳的办法是先在钢琴卷帘里听一遍实际音高，再决定移调。

\subsection{声部密度与留白}
\label{sec:density-space}

\term{声部密度}{voice-density}是同一时刻有多少条独立声部在说不同的话；\term{留白}{space}是故意不写满的空隙。两者是一枚硬币的两面：密度靠留白才听得清，留白靠密度才有对比。

三条经验：

\begin{itemize}
\item 一个功能层同时只留 1—2 条“说话”的声部；其余声部做长音、重复或退出。
\item 低频、中频、高频各要有主次；不要三层的重音都落在同一拍。
\item 副歌不是“所有轨都进来”，而是“每层都比主歌更完整”；留一条轨到桥段或尾奏再进，比一开始全上更有空间。
\end{itemize}

\subsection{对位与呼应}
\label{sec:counterpoint-call}

\term{对位}{counterpoint}是两条以上旋律线同时进行、各自独立又互相协调；\term{呼应句}{call-and-response}是一句先问、另一句后答。它们不是古典乐专属：流行歌的副旋律、爵士的问答、电子舞曲的 Lead 与 Bass 互动都用同一套听觉逻辑。

写对位先做三件事：

\begin{enumerate}
\item \textbf{错音区}：主旋律在中高音，副旋律就下到中低音；不要两条线在同一频段抢。
\item \textbf{错节奏}：主旋律走长句，副旋律就用短点或反拍；不要同时落在每个强拍。
\item \textbf{错方向}：一条上行，另一条优先下行或保持；同向平行容易糊成一条。
\end{enumerate}

呼应句则更简单：主句留一个“问号”（短停顿或未解决音），答句在同一音区或低一个八度接住。答句不一定要长，常常两三个音就够。

\subsection{织体类型}
\label{sec:texture-types}

\term{织体}{texture}是各声部在音区、节奏密度与音色上的组合方式。配器时常用的四类：

\begin{longtable}{p{0.18\textwidth}p{0.34\textwidth}p{0.38\textwidth}}
\toprule
织体 & 听感 & 典型用法 \\
\midrule
\endfirsthead
\toprule
织体 & 听感 & 典型用法 \\
\midrule
\endhead
主调织体 & 一条主旋律＋和弦伴奏，主次清楚 & 流行、摇滚、民谣的主歌 \\
复调织体 & 两条以上旋律同时独立进行 & 爵士问答、弦乐对位、影视主题对话 \\
织体分层 & 低频/主体/空气/节奏各占一层 & 电子舞曲、影视配乐、流行副歌 \\
固定音型 & 同一短句反复，变化靠位移/加花 & 摇滚 Riff、电子 Ostinato、国风循环 \\
\bottomrule
\end{longtable}

\section{八种曲风}
\label{sec:genres-eight}

每一节都按同一顺序：典型编制、节奏与律动、和声与旋律、段落结构、标志性手法、夜半骨架。所有轨数都是“教学起点”，不是必须；真正要守住的是功能分工。


\section{横向对照与练习}
\label{sec:genres-compare}

把八种曲风放在同一张表里，差别会更清楚：

\begin{longtable}{p{0.13\textwidth}p{0.25\textwidth}p{0.25\textwidth}p{0.25\textwidth}}
\toprule
曲风 & 时间骨架 & 和声/旋律重点 & 最容易犯的错 \\
\midrule
\endfirsthead
\toprule
曲风 & 时间骨架 & 和声/旋律重点 & 最容易犯的错 \\
\midrule
\endhead
流行 & 2/4 军鼓、1/3 底鼓 & 钩子清楚、句尾留白 & 副歌全轨同时推满 \\
摇滚 & 反拍清楚、Riff 驱动 & 强和弦、短句重复 & 吉他与贝斯完全同音 \\
电子舞曲 & 四拍四、8 小节块 & 循环动机、低频冲击 & Kick/Sub/Bass 三层抢低频 \\
嘻哈 & Boom-bap/Trap 口袋 & 采样循环、Hook 短句 & 808 与 Kick 同时满 \\
爵士 & Swing、Walking Bass & ii–V–I、Comping & 把 Swing 写成平均八分 \\
民谣 & 指弹固定音型、稀疏 & I–IV–V–vi、叙事 & 完整鼓组破坏呼吸 \\
影视配乐 & 画面驱动、Ostinato & 主题动机、调式色彩 & 全曲满编、没有留白 \\
国风 & 中速舒展、气口 & 五声音阶、装饰音 & 五声音阶贴标签、堆民乐 \\
\bottomrule
\end{longtable}

\begin{exercise}
选一个你不熟悉的曲风，按该节的“夜半骨架”做 8 小节：
\begin{enumerate}
\item 先建轨、命名、写段落标记；不要先找 Preset。
\item 只写节奏组与和声层，Solo 听一遍，确认低频不打架。
\item 再加旋律/人声与一个色彩层；每加一层，静音一次，确认它真的在帮忙。
\item 最后回答：如果去掉哪一条轨，音乐会变差？如果答不出，说明那条轨可能只是习惯。
\end{enumerate}
\end{exercise}

\section{本章要点}
\label{sec:genres-takeaways}

\begin{itemize}
\item 曲风是“编制—律动—和声—结构—标志手法”的组合，不是音色标签。
\item 音域、移调、密度、留白、对位、织体是八种风格共用的四把尺。
\item 每种风格先守功能分工，再谈乐器数量；少而清楚通常比多而糊更专业。
\item 夜半骨架要按“建轨—写骨架—加层—Solo 检查”的顺序落地，不要先堆 Preset。
\item 风格之间可以借手法，但不能借走对方的标志；借之前先问“它还认得出自己吗”。
\end{itemize}

\section{扩展阅读}
\label{sec:11-genres-further-reading}
\begin{itemize}
  \item \href{https://en.wikipedia.org/wiki/Music_genre}{Music genre（音乐曲风）}：看曲风怎样按乐器、织体与用法分类。
  \item \href{https://en.wikipedia.org/wiki/Arrangement}{Arrangement（编曲）}：回到“选择”本身，而不是背预设。
  \item \href{https://www.youtube.com/watch?v=M6WhsnWIhcM}{用编曲层与音色定义风格}：看不同层与音色怎样组合出一种风格。
  \item \href{https://www.youtube.com/watch?v=uO4_OuZHh1Q}{同一套鼓在不同曲风里的形态}：对照鼓组写法怎样随曲风改变。
\end{itemize}

\begin{helpmapping}
\begin{longtable}{p{0.45\textwidth}p{0.45\textwidth}}
\toprule
小节 & 帮助锚点 \\
\midrule
为什么把曲风独立成章 & \texttt{arrangement-genres-why} \\
音域与移调 & \texttt{arrangement-genres-range} \\
声部密度与留白 & \texttt{arrangement-genres-density} \\
对位与呼应 & \texttt{arrangement-genres-counterpoint} \\
织体类型 & \texttt{arrangement-genres-texture} \\
横向对照与练习 & \texttt{arrangement-genres-compare} \\
\bottomrule
\end{longtable}
\end{helpmapping}
